% !TeX root = ../../Book3.tex
\chapter{Riemann \texorpdfstring{\(\zeta\)}{\AnalysisPdfZeta} 函数}
\label{b3:ch:26}
\BookChapterNumber{b3:ch:26}
% 本章摘要（不计入编号）
\begin{chapterabstract}
Riemann zeta 函数最初由整数的 Dirichlet 级数定义，却通过 theta 的 Mellin 变换获得全平面延拓。本章证明其唯一极点、函数方程与平凡零点，明确尚待下一章处理的边界无零点问题；末节将非平凡零点编码在阶一整函数 \(\xi\) 的典范乘积中。
\end{chapterabstract}

% 本章目标（不计入编号）
\begin{learninggoals}
能够重建 theta 积分的分割与换元，核对 Gamma 因子和极点留数；区分原级数、延拓函数与完备函数的定义域；能准确说明零点对称性、已证零点范围及 Riemann 猜想的内容。
\end{learninggoals}

素数由 Euler 乘积进入 \(\zeta\)，全平面对称却来自高斯求和。把两个表达式连接起来，需要先给它们共同的 Mellin 积分，再把原点的两个显式奇性扣除。所需工具为\chapref{b3:ch:21}、\chapref{b3:ch:24}和\chapref{b3:ch:25}，末节另用 Hadamard 分解。Edwards 的《Riemann's Zeta Function》\cite{Edwards2001Riemann}可供继续阅读零点与显式公式的联系。

% !TeX root = ../../Book3.tex
\section{Dirichlet 级数与 Euler 乘积}
\label{b3:sec:2601}

本章把上一章的收敛级数扩展为全平面的亚纯函数。首先保留它与素数的直接关系；随后用高斯的参数反演得到在原级数收敛域之外仍有意义的表达式。

\subsection{在 \texorpdfstring{\(\operatorname{Re}s>1\)}{Re(s)>1} 上的收敛性}
\label{b3:subsec:260101}

\begin{definition}\label{b3:def:26-zeta}
在 \(\operatorname{Re}s>1\) 上定义\term[riemann-zeta-hanshu]{Riemann zeta 函数}[Riemann zeta function]
\(\zeta(s)\coloneqq\sum_{n\ge1}n^{-s}\)。
\end{definition}
\symidx{\(\zeta(s)\)：Riemann zeta 函数}
对紧集 \(\operatorname{Re}s\ge1+\delta\)，由 \(\sum n^{-1-\delta}<\infty\) 得局部一致绝对收敛，故全纯；任意固定次导数可逐项计算，因为 \((\log n)^k\) 被小幂 \(n^{\delta/2}\) 控制。在正实轴，\(\zeta(\sigma)>1\)，且 \(\zeta(\sigma)\to1\) 当 \(\sigma\to\infty\)。

\subsection{Euler 乘积}
\label{b3:subsec:260102}

\cref{b3:thm:25-euler-product}给 \(\zeta(s)=\prod_p(1-p^{-s})^{-1}\)。更强的是前章已经建立收敛全纯对数
\(\log\zeta(s)=\sum_{p,r\ge1}p^{-rs}/r\)，从而在这个半平面无零点。若把乘积形式延伸到 \(\operatorname{Re}s\le1\)，它通常不再收敛；后面的解析延拓不能被称为该无穷乘积在原意义下的值。

\ProseParagraphBreak

绝对收敛还给出不依赖虚部的上下界。由原级数及 Möbius 逆级数，
\[
 |\zeta(\sigma+\mathrm it)|\le\zeta(\sigma),\qquad
 \left|\frac1{\zeta(\sigma+\mathrm it)}\right|
 \le\sum_{n\ge1}\frac{|\mu(n)|}{n^\sigma}\le\zeta(\sigma),
 \quad\sigma>1.
\]
所以在每条固定的直线 \(\sigma>1\) 上，函数既有上界，也与零保持正距离。不过当 \(\sigma\downarrow1\) 时，这个下界趋于零，因而不能排除边界零点。无零点区域是否能够扩大，必须从新的信息判断。

\subsection{素数与 \texorpdfstring{\(\zeta\)}{\AnalysisPdfZeta} 函数}
\label{b3:subsec:260103}

\(\zeta\) 将所有整数编码，而其对数只在素数幂上有系数。沿实轴 \(\sigma\downarrow1\)，\(\zeta(\sigma)\) 的发散迫使素数无限多，也迫使 \(\sum_p1/p\) 发散。这类结论仍比素数定理弱：它们不能确定 \(\pi(x)\) 的渐近常数，后者需要控制函数在边界直线上的零点。

% !TeX root = ../../Book3.tex
\section{Mellin 变换与 \texorpdfstring{\(\theta\)}{\AnalysisPdfTheta} 函数}
\label{b3:sec:2602}

Gamma 因子出现的原因不是人为修饰：对一个高斯项做 Mellin 积分，二次尺度自然把参数变为 \(s/2\)，并产生 \(\pi^{-s/2}\Gamma(s/2)\)。

\subsection{Mellin 变换}
\label{b3:subsec:260201}

记 \(\vartheta(t)=\sum_{n\in\mathbb Z}e^{-\pi n^2t}\)、\(h(t)=(\vartheta(t)-1)/2=\sum_{n\ge1}e^{-\pi n^2t}\)，\(t>0\)。对 \(\sigma>1\) 可用 Tonelli 估计绝对值并换序：
\[
\int_0^\infty h(t)t^{s/2-1}\,\mathrm{d}t
=\sum_{n\ge1}(\pi n^2)^{-s/2}\Gamma(s/2)
=\pi^{-s/2}\Gamma(s/2)\zeta(s).
\]
绝对积分为 \(\pi^{-\sigma/2}\Gamma(\sigma/2)\zeta(\sigma)<\infty\)，所以复参数换序并非仅由每一项可积推出。

\subsection{Gamma 函数}
\label{b3:subsec:260202}

Gamma 函数的无零点性使这个乘法不会在右半平面创造新零点。其极点以后却会抵消 \(\zeta\) 的负偶数零点，所以除去 Gamma 因子时必须单独查看局部阶数。特别是 \(s=0\)，Gamma 有极点而 \(\zeta\) 最终并不为零，不能把所有 Gamma 极点一概称为平凡零点的位置。

\subsection{\texorpdfstring{\(\theta\)}{\AnalysisPdfTheta} 函数}
\label{b3:subsec:260203}

由\secref{b3:sec:2401}的高斯求和，\(\vartheta(t)=t^{-1/2}\vartheta(1/t)\)。于是
\[
h(t)=t^{-1/2}h(1/t)+\frac12(t^{-1/2}-1).
\]
这是延拓的关键等式：右侧第一项在 \(t\downarrow0\) 时带指数小因子，另外两项完整记录原点的幂次奇性。只减去常数一还不能使原点积分对所有参数收敛；还必须处理 \(t^{-1/2}\) 项。

\subsection{热核观点}
\label{b3:subsec:260204}

函数 \(\vartheta(t)\) 也可视为圆周高斯热核在原点的值，\(n^2\) 对应离散频率的平方。减去一是在去掉零频率。小时间的 \(t^{-1/2}\) 主项反映一维局部尺度；Mellin 积分把这个幂次转成 \(s=1\) 的极点。这个解释无需将谱语言当作证明，全部常数已由高斯换元确定。

\ProseParagraphBreak

为了使热核的归一化完全明确，在周期一的圆周上取 Fourier 模态 \(e^{2\pi\mathrm inx}\)。算子 \(-\partial_x^2\) 的对应特征值为 \(4\pi^2n^2\)，故热核在时间 \(\tau\) 的对角值是
\[
 K_\tau(0)=\sum_{n\in\mathbb Z}e^{-4\pi^2n^2\tau}
          =\vartheta(4\pi\tau).
\]
令 \(t=4\pi\tau\)，才得到前面使用的 \(\vartheta(t)\)。时间参数若不同时缩放，会使 Mellin 变换中的 \(\pi\) 因子发生变化。对角热核的大小与变换公式都可以从 Fourier 级数逐项验证，无需额外假定抽象谱定理。

\ProseParagraphBreak

小时间的余项也有直接界。对 \(u\ge1\)，利用 \(n^2\ge1+3(n-1)\) 可得
\[
 0<h(u)\le\frac{e^{-\pi u}}{1-e^{-3\pi u}}
          \le\frac{e^{-\pi u}}{1-e^{-3\pi}}.
\]
代入参数反演，便有
\[
 \vartheta(t)=t^{-1/2}
       +O\bigl(t^{-1/2}e^{-\pi/t}\bigr),\qquad0<t\le1.
\]
这比仅说余项趋于零强得多。对任意固定的复参数紧集，指数小量乘上 \(t\) 的任意固定负幂仍可积；因此扣去主项后能够穿过任意有限条竖直线。下一节的全平面延拓正是把这一估计写成一条统一的积分公式。

% !TeX root = ../../Book3.tex
\section{解析延拓}
\label{b3:sec:2603}

延拓的困难集中在积分的端点，而不在每个固定正实数处的被积函数。这里给出两条路线：Euler--Maclaurin 方法逐步减去端点主项，适合得到截断公式和误差；theta 方法利用反演，将两个端点合成一个指数收敛积分，同时显露 \(s\) 与 \(1-s\) 的对称性。希望得到具体截断误差时，可以选读第一小节的 Euler--Maclaurin 方法。

\OptionalSubsection{Euler–Maclaurin 方法}{b3:subsec:260301}

\eqref{b3:eq:25-zeta-integral}已在 \(\operatorname{Re}s>1\) 给出
\[
 \zeta(s)=\frac{s}{s-1}
      -s\int_1^\infty\{t\}t^{-s-1}\,\mathrm dt.
\]
因为 \(0\le\{t\}<1\)，右侧积分在 \(\operatorname{Re}s>0\) 仍绝对收敛。要继续向左，就不能只估计分数部分的大小，还要利用它减去平均值以后每个周期的积分为零。一个均值为零的周期函数有有界周期原函数；分部积分把这个原函数留下，同时让权重多衰减一个 \(t^{-1}\)。反复实行这一操作，得到\term[euler-maclaurin-qiuhe-gongshi]{Euler--Maclaurin 求和公式}[Euler--Maclaurin summation formula]。

\ProseParagraphBreak

为使每次积分中的常数唯一确定，先定义\term[bernoulli-duoxiangshi]{Bernoulli 多项式}[Bernoulli polynomial]。置 \(B_0(x)=1\)；对整数 \(j\ge1\)，以
\[
 B_j'(x)=jB_{j-1}(x),
 \quad \int_0^1B_j(x)\,\mathrm dx=0
\]
唯一确定多项式 \(B_j\)。导数条件留下一个常数，零平均条件恰好确定它。前几项为
\[
 \begin{aligned}
 B_1(x)&=x-\frac12,\quad
 B_2(x)=x^2-x+\frac16,\\
 B_3(x)&=x^3-\frac32x^2+\frac12x,\\
 B_4(x)&=x^4-2x^3+x^2-\frac1{30}.
 \end{aligned}
\]
数 \(B_j\coloneqq B_j(0)\) 称为\term[bernoulli-shu]{Bernoulli 数}[Bernoulli number]。省略自变量的 \(B_j\) 是常数，带自变量的 \(B_j(x)\) 是多项式；例如 \(B_1=-1/2\)、\(B_2=1/6\)、\(B_4=-1/30\)。这里固定了 \(B_1\) 的负号约定。

\ProseParagraphBreak

记 \(\widetilde B_j(t)=B_j(\{t\})\)。这些函数有界且一周期；对 \(j\ge2\)，它们在整数点也是连续的，因为
\[
 B_j(1)-B_j(0)=j\int_0^1B_{j-1}(x)\,\mathrm dx=0.
\]
此外，递推中的导数与零平均条件给出反射关系
\[
 B_j(1-x)=(-1)^jB_j(x).
\]
具体地，若此式对 \(j-1\) 成立，则 \((-1)^jB_j(1-x)\) 与 \(B_j(x)\) 有相同导数和相同的零积分；它们的差为积分也为零的常数，故相等。\(j=0\) 提供归纳起点。结合整数端点相等性可知，奇数 \(j\ge3\) 时两端值既相等又互为相反数，因而均为零。这正是截断公式只留下偶数 Bernoulli 项的原因。

\begin{proposition}[Euler--Maclaurin 截断公式]\label{b3:prop:26-euler-maclaurin}
设 \(N,m\ge1\) 为整数，Bernoulli 多项式由 \(B_0(x)=1\)、\(B_j'=jB_{j-1}\) 及 \(\int_0^1B_j=0\)（\(j\ge1\)）确定，记 \(B_j=B_j(0)\)、\(\widetilde B_j(t)=B_j(\{t\})\)，并令
\[
 (s)_0=1,\quad (s)_j=s(s+1)\cdots(s+j-1)\quad(j\ge1).
\]
原在 \(\operatorname{Re}s>1\) 由 \(\sum_{n\ge1}n^{-s}\) 定义的 \(\zeta\) 有亚纯延拓到 \(\operatorname{Re}s>1-2m\)，且在该半平面有
\begin{equation}\label{b3:eq:26-euler-maclaurin}
 \begin{aligned}
 \zeta(s)={}&\sum_{n<N}n^{-s}+\frac{N^{1-s}}{s-1}+\frac{N^{-s}}2\\
 &+\sum_{r=1}^m\frac{B_{2r}}{(2r)!}
       (s)_{2r-1}N^{1-s-2r}+R_m(s,N),
 \end{aligned}
\end{equation}
其中正实数的复幂用实对数定义，\(s=1\) 处按亚纯函数解释，并且
\[
 R_m(s,N)=-\frac{(s)_{2m}}{(2m)!}
       \int_N^\infty\widetilde B_{2m}(t)t^{-s-2m}\,\mathrm dt.
\]
余项在所述半平面全纯；若 \(\sigma=\operatorname{Re}s>1-2m\)，则
\begin{equation}\label{b3:eq:26-euler-maclaurin-bound}
 |R_m(s,N)|\le
 \frac{|(s)_{2m}|}{(2m)!}\cdot
 \frac{\|B_{2m}\|_{L^\infty(0,1)}}{\sigma+2m-1}
 N^{1-\sigma-2m}.
\end{equation}
不同 \(m,N\) 所得的延拓在共同定义域上一致。
\end{proposition}
\begin{proof}
首先限制 \(\sigma>1\)，保证最初的无穷求和与各个端点极限合法。由\cref{b3:prop:25-integral-summation}，对整数 \(N\) 后面的尾项有
\[
 \sum_{n>N}n^{-s}
 =-N^{1-s}+s\int_N^\infty\lfloor t\rfloor t^{-s-1}\,\mathrm dt.
\]
加入 \(N^{-s}\)，再代入 \(\lfloor t\rfloor=t-1/2-\widetilde B_1(t)\)，计算前两项幂积分，得到
\[
 \sum_{n\ge N}n^{-s}
 =\frac{N^{1-s}}{s-1}+\frac12N^{-s}
       -s\int_N^\infty\widetilde B_1(t)t^{-s-1}\,\mathrm dt.
\]
虽然 \(\widetilde B_1\) 在整数点跳跃，但它在每个开单位区间内为 \(\widetilde B_2/2\) 的导数。逐区间分部积分时，\(B_2(1)=B_2(0)\) 使内部端点相消；其在无穷端的有界性与权重衰减使该端点消失。下端点给出
\[
 -s\int_N^\infty\widetilde B_1(t)t^{-s-1}\,\mathrm dt
 =\frac{B_2}{2!}sN^{-s-1}
       -\frac{(s)_2}{2!}\int_N^\infty
                     \widetilde B_2(t)t^{-s-2}\,\mathrm dt.
\]
这是 \(m=1\) 的公式，也核定了余项前的负号。

假设已经得到偶数阶 \(2r\) 的余项。再分部积分时，原函数为 \(\widetilde B_{2r+1}/(2r+1)\)，其两个整数端点都为零，因此没有新的边界贡献；随后使用 \(\widetilde B_{2r+2}/(2r+2)\) 再分部积分，下端点恰给下一项
\[
 \frac{B_{2r+2}}{(2r+2)!}(s)_{2r+1}N^{-s-2r-1}.
\]
余项仍有负号，阶数变为 \(2r+2\)。内部整数端点依靠相邻两端值相等而相消。在 \(\sigma>1\) 中所有无穷端点均为零，故归纳得到所述有限展开。

最后利用 \(\widetilde B_{2m}\) 有界，对 \(\sigma>1-2m\) 计算
\[
 \int_N^\infty t^{-\sigma-2m}\,\mathrm dt
 =\frac{N^{1-\sigma-2m}}{\sigma+2m-1}.
\]
这证明余项界。在该半平面任一紧集上，指数离临界值有正距离；参数微分增加的 \(\log t\) 幂仍可积。因此余项全纯，右端只可能在一处有极点，并在原半平面与原级数相同，故给出延拓。任意两种截断表达式均在 \(\sigma>1\) 相同，恒等定理保证它们在连通的共同半平面相同。
\end{proof}

一个实际计算应同时报告近似值和误差。取 \(m=1\)，因为 \(B_2(x)=x^2-x+1/6\) 在 \([0,1]\) 上的最大绝对值为 \(1/6\)，有
\[
 |R_1(s,N)|\le
   \frac{|s(s+1)|}{12(\sigma+1)}N^{-\sigma-1},
 \quad \sigma>-1.
\]
例如 \(s=2\)、\(N=10\) 时，
\[
 \left|\zeta(2)-
   \left(\sum_{n=1}^{9}\frac1{n^2}+\frac1{10}
                   +\frac1{200}+\frac1{6000}\right)\right|
 \le\frac1{6000}.
\]
右端是已知数，左端只需计算九个有理数的和。若 \(|\operatorname{Im}s|\) 增大，分子 \(|s(s+1)|\) 也会增大；固定 \(N\) 的这条误差估计并不在无界竖线上统一变小。

\ProseParagraphBreak

截断公式也能计算原级数发散处的延拓值。取 \(N=1\)、\(m=1\)、\(s=0\)，所有含 \(s\) 的项为零，得到 \(\zeta(0)=-1/2\)。若求 \(s=-1\)，则不能把它代入要求 \(\operatorname{Re}s>-1\) 的同一版本；须取 \(m=2\)，此时合法域为 \(\operatorname{Re}s>-3\)。因为 \((s)_3\) 与 \((s)_4\) 在 \(-1\) 都为零，剩下
\[
 \zeta(-1)=-\frac12+\frac12-\frac1{12}=-\frac1{12}.
\]
这些都是亚纯延拓的取值，不是发散正项级数的通常和。下面的主线将在不调用本小节的情况下建立同一延拓，并由其对称性重新得到相应特殊值。

\subsection{\texorpdfstring{\(\theta\)}{\AnalysisPdfTheta} 方法}
\label{b3:subsec:260302}

回到\cref{b3:prop:26-theta-mellin}。先只取 \(\operatorname{Re}s>1\)，使原 Mellin 积分绝对收敛。在 \(t=1\) 处分成两段，对 \((0,1)\) 上的 \(h(t)\) 使用\eqref{b3:eq:26-theta-inversion}：
\[
 \begin{aligned}
 \int_0^1h(t)t^{s/2-1}\,\mathrm dt
 ={}&\int_0^1 t^{-1/2}h(1/t)t^{s/2-1}\,\mathrm dt\\
 &+\frac12\int_0^1
       \left(t^{-1/2}-1\right)t^{s/2-1}\,\mathrm dt.
 \end{aligned}
\]
两个显式幂积分可分别计算，因为 \(\operatorname{Re}s>1\) 同时保证它们在原点可积，结果为
\[
 \frac12\left(\frac2{s-1}-\frac2s\right)
 =\frac1{s-1}-\frac1s.
\]
这里负号来自被减去的常数项；不能为了得到一个对称公式而省略符号核算。

\ProseParagraphBreak

对另一项令 \(u=1/t\)。此时 \(\mathrm dt=-u^{-2}\,\mathrm du\)，上下限由 \((0,1)\) 变为 \((\infty,1)\)。反转上下限以后，
\[
 \begin{aligned}
 \int_0^1t^{-1/2}h(1/t)t^{s/2-1}\,\mathrm dt
 &=\int_1^\infty h(u)u^{-(s+1)/2}\,\mathrm du\\
 &=\int_1^\infty h(u)u^{(1-s)/2-1}\,\mathrm du.
 \end{aligned}
\]
原先 \((1,\infty)\) 上的积分仍带着幂次 \(s/2-1\)。将两部分相加得到
\begin{equation}\label{b3:eq:26-completed-integral}
 \begin{aligned}
 \pi^{-s/2}\Gamma(s/2)\zeta(s)
 ={}&-\frac1s-\frac1{1-s}\\
 &+\int_1^\infty h(t)
      \left(t^{s/2-1}+t^{(1-s)/2-1}\right)\,\mathrm dt.
 \end{aligned}
\end{equation}
这时才出现 \(s\) 与 \(1-s\) 的一对指数。选在一处分割，是因为一是变换 \(t\mapsto1/t\) 的固定点；它让两段最后具有同一个积分区间。

\ProseParagraphBreak

现在检查右端的适用域。记其中的积分为 \(I(s)\)。给定任意紧集 \(K\subset\mathbb C\)，\(s\) 的实部在 \(K\) 上有上下界，故存在实数 \(A\)，使所有 \(s\in K\)、\(t\ge1\) 满足
\[
 \left|t^{s/2-1}\right|
 +\left|t^{(1-s)/2-1}\right|\le2t^A.
\]
由\cref{b3:prop:26-theta-endpoints}，积分核由 \(Ce^{-\pi t}t^A\) 控制。第 \(k\) 次参数导数至多再添 \(2^{-k}(\log t)^k\)；指数衰减仍保证可积。因此 \(I(s)\) 在整个复平面全纯，且可任意次在积分号下求导。特别地，
\[
 I'(s)=\frac12\int_1^\infty h(t)\log t
       \left(t^{s/2-1}-t^{(1-s)/2-1}\right)\,\mathrm dt.
\]
同一个支配还给数值截断的尾界：在 \(K\) 上把积分截在 \(T\ge1\) 时，余量不超过常数倍 \(\int_T^\infty e^{-\pi t}t^A\,\mathrm dt\)；求导后的余量只增加已知的对数幂。因此这是一条真正收敛的新表示，而不是对原发散级数作形式操作。

\subsection{亚纯延拓}
\label{b3:subsec:260303}

\begin{theorem}[Riemann zeta 函数的亚纯延拓]\label{b3:thm:26-continuation}
在 \(\operatorname{Re}s>1\) 上由绝对收敛级数
\(\zeta(s)=\sum_{n=1}^\infty e^{-s\log n}\) 定义的 Riemann zeta 函数，唯一亚纯延拓到全复平面。该延拓唯一的极点是 \(s=1\) 处的单极点，其留数为一；在 \(s=0\) 处全纯且 \(\zeta(0)=-1/2\)。
\end{theorem}
\begin{proof}
用 \(h(t)=\sum_{n\ge1}e^{-\pi n^2t}\) 定义
\[
 F(s)=-\frac1s-\frac1{1-s}+I(s),
 \quad
 I(s)=\int_1^\infty h(t)
      \left(t^{s/2-1}+t^{(1-s)/2-1}\right)\,\mathrm dt.
\]
上小节已经证明 \(I\) 为整函数。因此 \(F\) 恰在零、一两处有单极点，留数分别为 \(-1\)、\(1\)。定义亚纯函数
\[
 Z(s)=\frac{\pi^{s/2}}{\Gamma(s/2)}F(s).
\]
\cref{b3:thm:21-product}保证前面的乘子在全平面整，故 \(Z\) 不会在其他点产生新极点。在原点，用\eqref{b3:eq:26-reciprocal-gamma-local}可得
\[
 Z(s)=\left(\frac s2+O(s^2)\right)
          \left(-\frac1s+O(1)\right)
      =-\frac12+O(s).
\]
故该点可去，延拓值为 \(-1/2\)。在 \(s=1\) 处，乘子全纯且值为 \(\pi^{1/2}/\Gamma(1/2)=1\)，故 \(Z\) 仍有单极点，留数为一，不能被抵消。

由\eqref{b3:eq:26-completed-integral}，\(Z(s)=\zeta(s)\) 在 \(\operatorname{Re}s>1\) 成立，因而 \(Z\) 给出所求延拓。若还有另一个全平面亚纯延拓，两者在这个非空开集相同；亚纯函数恒等定理使它们在连通域 \(\mathbb C\) 上处处相同。这证明唯一性。
\end{proof}

以后仍将这个唯一延拓记作 \(\zeta\)。特别是谈到 \(\zeta(0)\) 或 \(\zeta(-1)\) 时，说的是该亚纯函数的有限值，原级数在这些点并不收敛。

\ProseParagraphBreak

在负偶数 \(-2m\)、\(m\ge1\) 处，刚才的 \(F\) 全纯，而 \(1/\Gamma(s/2)\) 有单零点，因此目前已经知道 \(\zeta(-2m)=0\)。但“至少有一个零因子”还不能确定所得零点恰为一阶：还需排除 \(F(-2m)=0\)。\secref{b3:sec:2604}的对称关系与\secref{b3:sec:2605}的零点论证将补上这一点。对 \(s=0\) 则必须单独处理，因为 \(F\) 本身在那里有极点；上面的局部计算说明它与负偶数的情形不同。

\subsection{\texorpdfstring{\(s=1\)}{s=1} 处的极点}
\label{b3:subsec:260304}

留数给出最强的发散项，还没有给出减去它以后留下的常数。这里无需从复杂的 theta 积分中辨认常数，而应回到能直接看见调和和的部分求和公式。

\begin{proposition}[一处的 Laurent 常数]\label{b3:prop:26-laurent-at-one}
令 \(\zeta\) 为 \(\sum_{n\ge1}n^{-s}\) 在 \(\operatorname{Re}s>1\) 所定义函数的唯一亚纯延拓，令
\[
 \gamma=\lim_{N\to\infty}\left(\sum_{n=1}^N\frac1n-\log N\right)
\]
为 Euler 常数。则 \(\zeta(s)-1/(s-1)\) 在 \(s=1\) 可去，且其延拓值为 \(\gamma\)。因此在该点邻域有收敛的 Laurent 展开，其前两项为
\begin{equation}\label{b3:eq:26-laurent-at-one}
 \zeta(s)=\frac1{s-1}+\gamma+O(s-1).
\end{equation}
\end{proposition}
\begin{proof}
由\eqref{b3:eq:25-zeta-integral}，在 \(\operatorname{Re}s>1\) 时，
\[
 \zeta(s)-\frac1{s-1}
 =1-s\int_1^\infty\{t\}t^{-s-1}\,\mathrm dt.
\]
右侧在 \(\operatorname{Re}s>0\) 全纯：对任意紧集，可取 \(\delta>0\) 使 \(\operatorname{Re}s\ge\delta\)，被积函数及其各阶参数导数分别被常数倍 \(t^{-1-\delta}(\log t)^k\) 控制。由恒等定理，这个表达式与已经构造的 \(\zeta(s)-1/(s-1)\) 在其共同域一致。所以在 \(s=1\) 可直接代入，得到常数
\[
 1-\int_1^\infty\frac{\{t\}}{t^2}\,\mathrm dt.
\]

为了识别它，先取整数上限 \(N\ge2\)。在 \([k,k+1)\) 上，\(\{t\}=t-k\)，故
\[
 \int_k^{k+1}\frac{\{t\}}{t^2}\,\mathrm dt
 =\int_k^{k+1}\left(\frac1t-\frac{k}{t^2}\right)\,\mathrm dt
 =\log\frac{k+1}{k}-\frac1{k+1}.
\]
从 \(k=1\) 加到 \(N-1\)，对数项望远镜相消，得到
\[
 \int_1^N\frac{\{t\}}{t^2}\,\mathrm dt
 =\log N-H_N+1,
 \quad H_N=\sum_{n=1}^N\frac1n.
\]
令 \(N\to\infty\)，利用\cref{b3:prop:21-euler-constant}已证明的极限 \(H_N-\log N\to\gamma\)，便得所需常数为 \(\gamma\)。最后，去掉主部后的函数在一处全纯，故其 Taylor 展开确实给出 \(\gamma+O(s-1)\)，而不只是沿实轴的一次极限。
\end{proof}

对实数 \(s\downarrow1\)，这说明 \(\zeta(s)\) 的发散速度是 \(1/(s-1)\)，减去该主项以后趋于 \(\gamma\)。同一个常数也曾出现在 Gamma 的无穷乘积中；两个位置的联系都通过 \(H_N-\log N\) 的极限得到，因而常数相同有明确依据。接下来的函数方程计算将在 \(s=0\) 附近使用\eqref{b3:eq:26-laurent-at-one}，那里需要同时保留主部与常数项，才能看见各因子的一阶项如何抵消。

% !TeX root = ../../Book3.tex
\section{函数方程}
\label{b3:sec:2604}

延拓积分的两个幂次在 \(s\) 与 \(1-s\) 互换时交换位置。这是函数方程的来源，而不是由零点图形推测出的对称。

\subsection{\texorpdfstring{\(\theta\)}{\AnalysisPdfTheta} 变换}
\label{b3:subsec:260401}

在\eqref{b3:eq:26-completed-integral}中，显式有理项和积分都在 \(s\mapsto1-s\) 下不变。所有换元都最初在绝对收敛域完成，而最终的等式是全平面的亚纯恒等式；不存在要求 \(s\) 与 \(1-s\) 同时落在原 Dirichlet 收敛域的矛盾。

\subsection{完备 \texorpdfstring{\(\zeta\)}{\AnalysisPdfZeta} 函数}
\label{b3:subsec:260402}

定义\term[wanbei-zeta-hanshu]{完备 zeta 函数}[completed zeta function]
\(\Lambda_\zeta(s)\coloneqq\pi^{-s/2}\Gamma(s/2)\zeta(s)\)。下标用来区别算术函数 \(\Lambda(n)\)。它只在零与一有单极点，留数分别为负一与一，并满足
\begin{equation}\label{b3:eq:26-functional-equation}
\Lambda_\zeta(s)=\Lambda_\zeta(1-s).
\end{equation}
\symidx{\(\Lambda_\zeta(s)\)：完备 zeta 函数}

\subsection{函数方程}
\label{b3:subsec:260403}

利用 Gamma 的反射和倍增公式，将\eqref{b3:eq:26-functional-equation}中的 Gamma 比值化简，得到另一形式
\[
\zeta(s)=2^s\pi^{s-1}\sin(\pi s/2)\Gamma(1-s)\zeta(1-s).
\]
化简可从 \(\Gamma(s/2)\Gamma(1-s/2)=\pi/\sin(\pi s/2)\) 和
\(\Gamma((1-s)/2)\Gamma(1-s/2)=2^s\sqrt\pi\Gamma(1-s)\) 出发。式中各因子在某些整数处分别无穷或为零，等式应按亚纯延拓解释，不把那些点的形式乘积直接当作普通数值计算。

\subsection{对称性}
\label{b3:subsec:260404}

在讨论零点以前，先以两个局部计算检验这些因子的配合。令 \(s\to0\)，由已知 Laurent 展开和 Gamma 在一处的导数，
\[
 \zeta(1-s)=-\frac1s+\gamma+O(s),\qquad
 \Gamma(1-s)=1+\gamma s+O(s^2).
\]
两者乘积的常数项恰好抵消。再展开
\[
 2^s\pi^{s-1}=\pi^{-1}\bigl(1+s\log(2\pi)+O(s^2)\bigr),
 \qquad \sin(\pi s/2)=\frac{\pi s}2+O(s^3),
\]
得到
\[
 \zeta(s)=-\frac12-\frac12\log(2\pi)s+O(s^2).
\]
所以 \(\zeta'(0)=-\tfrac12\log(2\pi)\)。这里有限结果来自正弦零点与 \(\zeta(1-s)\) 极点的抵消，不是把两个无定义的点值直接相乘。

\ProseParagraphBreak

另一方面，对整数 \(m\ge1\)，在 \(s=1-2m\) 处各因子都是普通有限值，因而
\[
 \zeta(1-2m)=(-1)^m\frac{2(2m-1)!}{(2\pi)^{2m}}\zeta(2m).
\]
例如 \(\zeta(-1)=-1/12\)。负奇数值由正偶数值确定，而负偶数处由正弦因子产生零点，两种情形不可混淆。以上数值是延拓函数的值，不是把原发散 Dirichlet 级数作通常求和得到的值。

\ProseParagraphBreak

因为原级数实系数，\(\zeta(\overline s)=\overline{\zeta(s)}\)，延拓保持此式。完备函数还有 \(s\mapsto1-s\) 对称，所以非实非平凡零点按 \(\rho,\overline\rho,1-\rho,1-\overline\rho\) 出现；四点可能重合为两点，例如落在中线时。这个对称关系限制零点位置，却不强迫每条对称轨道都缩在中线上。

% !TeX root = ../../Book3.tex
\section{零点}
\label{b3:sec:2605}

函数方程与 Euler 乘积共同限制零点，但边界直线上的无零点性需要额外论证。本节先精确区分已经排除的区域与下一章要处理的边界。

\subsection{平凡零点}
\label{b3:subsec:260501}

\begin{proposition}\label{b3:prop:26-trivial-zeros}
在 \(\operatorname{Re}s<0\) 中，\(\zeta\) 的零点恰为负偶整数，且都是单零点；此外 \(\zeta(0)=-1/2\)。
\end{proposition}
\begin{proof}
在该半平面，\(1-s\) 实部大于一，故 \(\zeta(1-s)\ne0\)，\(\Gamma(1-s)\) 全纯非零。函数方程中只有正弦因子可能为零，其零点为负偶整数，且均单。零处的值则直接由延拓积分的 \(-1/s\) 项乘以 \(s/2+O(s^2)\) 得到。
\end{proof}
这些零点称为\term[pingfan-lingdian]{平凡零点}[trivial zero]；“平凡”指其位置直接由显式因子确定，不表示相关极限可省略计算。

\subsection{临界带}
\label{b3:subsec:260502}

除负偶数外，其他零点只能位于闭条带 \(0\le\operatorname{Re}s\le1\)，因为右半平面一以右由 Euler 乘积排除。开条带 \(0<\operatorname{Re}s<1\) 称为\term[linjie-dai]{临界带}[critical strip]。\chapref{b3:ch:27}将证明右边界一上无零，再由函数方程排除左边界上非零虚部的零点。当前使用“闭条带”保留了这一尚需证明的步骤。

\subsection{临界线}
\label{b3:subsec:260503}

直线 \(\operatorname{Re}s=1/2\) 称为\term[linjie-xian]{临界线}[critical line]。在这条线上，消去极点后的对称函数将取实值；这使寻找中线零点可以化成实函数的变号问题。但是实值性不保证函数必变号，也不排除带内其他零点。

\subsection{非平凡零点}
\label{b3:subsec:260504}

不属于负偶数的零点称为\term[feipingfan-lingdian]{非平凡零点}[nontrivial zero]。它们的重数按全纯函数局部零点阶数定义；关于它们是否全部为单零点，本书不作假设。实区间 \((0,1)\) 上没有零点：交错级数
\(\eta(s)=\sum(-1)^{n-1}n^{-s}\) 在 \(\operatorname{Re}s>0\) 全纯，并满足 \(\eta=(1-2^{1-s})\zeta\)。对实 \(0<s<1\)，将相邻两项配对得到严格正数，分母 \(1-2^{1-s}<0\)，故 \(\zeta(s)<0\)。严格正性可由第一对正项及其余非负对看出。

\subsection{Riemann 猜想}
\label{b3:subsec:260505}

交错级数还说明一个表达式的分母零点未必是延拓函数的极点。在 \(\operatorname{Re}s>0\) 中，恒等式
\[
 \eta(s)=(1-2^{1-s})\zeta(s)
\]
先在绝对收敛半平面通过偶数项分离得到，再按恒等定理延拓。因子 \(1-2^{1-s}\) 在
\(s_k=1+2\pi\mathrm ik/\log2\) 处有单零点。若 \(k\ne0\)，\(\zeta\) 在该点全纯，所以 \(\eta(s_k)=0\)，商式中的表面奇点可以去除；只有 \(k=0\) 对应真正的极点，此时 \(\eta(1)=\log2\ne0\)。因此这个商可以用于跨过直线一，但必须处理这些离散的表面奇点，不能把分母的每个零点都报告为 zeta 的极点。

\ProseParagraphBreak

对于实 \(0<s<1\)，也可从相邻配对获得定量符号：第一对为 \(1-2^{-s}>0\)，其余各对为非负数，故
\[
 \eta(s)\ge1-2^{-s}>0,
 \qquad
 \zeta(s)\le-\frac{1-2^{-s}}{2^{1-s}-1}<0.
\]
这个不等式只在实区间使用了项的正性；复参数时各项没有这样的序关系。因此不能把实轴无零点论证直接平移成临界带内的无零点论证。

\ProseParagraphBreak

\begin{conjecture}[Riemann]\label{b3:conj:26-rh}
Riemann zeta 函数的每个非平凡零点都满足 \(\operatorname{Re}s=1/2\)。
\end{conjecture}
这称为\term[riemann-caixiang]{Riemann 猜想}[Riemann hypothesis]。它比素数定理所需的边界无零点性强得多。有限高度的数值验证只约束被检查的有限区域，即使精度和零点计数都严格，也不覆盖无穷高处的零点。本章随后证明零点整体上无限多，这同样不能替代它们全部位于中线的断言。

% !TeX root = ../../Book3.tex
\OptionalSection{\texorpdfstring{\(\xi\)}{\AnalysisPdfXi} 函数与 Hadamard 乘积}{b3:sec:2606}

Gamma 因子消去平凡零点后，只剩两个显式极点。再乘一个二次多项式便得到整函数，因而可以应用整函数增长与典范乘积理论。

\subsection{Riemann \texorpdfstring{\(\xi\)}{\AnalysisPdfXi} 函数}
\label{b3:subsec:260601}

定义\term[riemann-xi-hanshu]{Riemann xi 函数}[Riemann xi function]
\[
\xi(s)\coloneqq\frac12s(s-1)\Lambda_\zeta(s).
\]
\symidx{\(\xi(s)\)：Riemann xi 函数}
它是整函数，\(\xi(0)=\xi(1)=1/2\)，满足 \(\xi(s)=\xi(1-s)\) 和共轭对称。其零点恰为非平凡零点，重数相同。负偶数处的 Gamma 单极点与 zeta 单零点抵消，并留下非零值，故这些点不在 \(\xi\) 的零点集合内。

\subsection{整函数的阶}
\label{b3:subsec:260602}

\begin{proposition}\label{b3:prop:26-xi-order}
整函数 \(\xi\) 的增长阶恰为一，且有无穷多个零点。
\end{proposition}
\begin{proof}
在 \(|s|\le r\) 上，延拓积分给
\[
 |\xi(s)|\le C(1+r)^2\left[1+
    \int_1^\infty e^{-\pi t}\bigl(t^{(r+1)/2}+1\bigr)\,\mathrm{d}t\right].
\]
右侧由 Gamma 积分控制，其对数为 \(O(r\log(r+2))\)，由正实轴 Stirling 公式得到，所以阶至多一。沿实轴 \(s\to+\infty\)，\(\zeta(s)\to1\)，Stirling 给
\[
 \log\xi(s)=\frac{s}{2}\log s+O(s),
\]
故阶至少一。
若只有有限零点，\cref{b3:thm:20-hadamard}使 \(\xi(s)=P(s)e^{as+b}\)，沿正实轴对数模仅为 \(O(s+\log s)\)，与上一渐近矛盾。因此零点无限多。
\end{proof}

\subsection{Hadamard 乘积}
\label{b3:subsec:260603}

增长上界与正实轴下界在这里承担不同作用。上界保证可以使用阶一的典范因子，并控制零点总数；下界排除有限零点乘一个线性指数的可能。仅有上界不能证明零点无限多，例如 \(e^s\) 也是阶一整函数，却没有零点。由于本章已经排除实区间内的零点，而平凡零点又被 Gamma 因子抵消，所得无穷多个零点确实都是非实的非平凡零点。

\ProseParagraphBreak

将零点按重数列为 \(\rho\)，Hadamard 分解给
\[
\xi(s)=e^{A+Bs}\prod_\rho(1-s/\rho)e^{s/\rho}.
\]
零点不存在于原点，故没有额外的 \(s^m\) 因子。收敛意义为属一典范乘积的正规收敛；\(\sum|\rho|^{-2}<\infty\) 由阶一零点估计得到。若要去掉指数修正，必须规定对称配对并另外证明配对和收敛，不应直接删除。

\subsection{零点分布}
\label{b3:subsec:260604}

乘积中的两个常数也有局部意义。取 \(s=0\)，每个典范因子等于一，因此可取 \(e^A=1/2\)。由于每个因子的对数导数在零点为零，\(B=\xi'(0)/\xi(0)\)。利用 Gamma 在零处的对数导数
\(\Gamma'(z)/\Gamma(z)=-1/z-\gamma+O(z)\)，以及上一节求得的 \(\zeta'(0)\)，有
\[
 \frac{\xi'(s)}{\xi(s)}
 =\frac1s+\frac1{s-1}-\frac12\log\pi
   +\frac12\frac{\Gamma'(s/2)}{\Gamma(s/2)}
   +\frac{\zeta'(s)}{\zeta(s)}.
\]
其中 \(1/s\) 与 Gamma 项的主部相消。取极限，得到
\[
 B=\log2+\frac12\log\pi-1-\frac\gamma2.
\]
这一步只确定乘积的指数因子，不提供单个零点的位置。它也是检验 Gamma 因子、\(1/2\) 常数和极点消去是否一致的一种方法。

\ProseParagraphBreak

Jensen 公式结合 \(\log M(r,\xi)=O(r\log(r+2))\) 给 \(N(r)=O(r\log(r+2))\)，其中 \(N(r)\) 为圆盘内零点总数。具体取外半径 \(2r\)，每个 \(|\rho|\le r\) 在 Jensen 和中贡献至少 \(\log2\)，故 \(N(r)\log2\le\log M(2r)-\log|\xi(0)|\)。这只是上界；精确的 Riemann--von Mangoldt 渐近还需要控制辐角，本章不由增长上界宣称该渐近已经成立。

\begin{proposition}[关于中点配对的乘积]
\label{b3:prop:26-centered-product}
置 \(F(z)=\xi(1/2+z)\)。从每一对非零零点 \(\alpha,-\alpha\) 中取一项，并保留重数，则
\[
 F(z)=F(0)\prod_{\{\alpha,-\alpha\}}
           \left(1-\frac{z^2}{\alpha^2}\right).
\]
该乘积在紧集上正规收敛，且不依赖各对之间的排列。
\end{proposition}
\begin{proof}
函数方程给 \(F(-z)=F(z)\)，故零点按相同重数成对。又 \(\zeta(1/2)<0\)、\(\Gamma(1/4)>0\) 且 \((1/2)(-1/2)<0\)，所以 \(F(0)>0\)，没有中心零点。平移不改变增长阶，对 \(F\) 使用 Hadamard 分解。每一对属一因子的乘积为
\[
 E_1(z/\alpha)E_1(-z/\alpha)=1-z^2/\alpha^2.
\]
由零点计数上界，\(\sum|\alpha|^{-2}<\infty\)，所以配对乘积正规收敛，并且可以任意重排各对。于是 \(F(z)=e^{a+bz}P(z)\)，其中 \(P\) 为偶函数且 \(P(0)=1\)。在原点附近 \(P\) 非零，\(F\) 的偶性给 \(e^{bz}=e^{-bz}\)。对此在零点求导得 \(b=0\)，再取值确定 \(e^a=F(0)\)。
\end{proof}

这种配对是在平移后的中点坐标中进行的，不能把它误读成原变量中任意顺序的 \(\prod_\rho(1-s/\rho)\)。配对收敛所依赖的是平方倒数可和，而单个倒数的绝对可和并未建立。乘积看起来更短，实际增加的是明确的对称分组规则。

\subsection{显式公式观点}
\label{b3:subsec:260605}

对典范乘积在避开零点的紧集取对数导数，有
\[
\frac{\xi'(s)}{\xi(s)}=B+\sum_\rho\left(\frac1{s-\rho}+\frac1\rho\right).
\]
每个括号为 \(O_K(|\rho|^{-2})\)，故正规收敛。另一方面，定义式将它写成 \(\zeta'/\zeta\)、Gamma 对数导数与有理项之和。素数幂级数与零点部分分式由此描述同一个亚纯函数。在反演围道中，跨越零点会取到由其重数决定的留数，这便是显式公式的基本机制；具体无穷零点求和还必须指定截断和误差，不能只把有限留数定理中的省略号延长到无穷。

% 本章小结（不计入编号）
\begin{chaptersummary}
theta 反演产生 \(s\) 与 \(1-s\) 的函数方程，Gamma 因子同时记录平凡零点的抵消。Euler 乘积排除右半平面一以右的零点；临界边界的排除还需要正性论证。\(\xi\) 将余下零点纳入整函数理论，增长上界与典范乘积提供了零点计数和部分分式的可控形式。
\end{chaptersummary}


\BookExercises

% 八题自拟：局部值、抵消、对称、Mellin缩放及增长估计。
\begin{exercise}\label{b3:ex:26-1}
已知 \(\zeta(4)=\pi^4/90\)。求 \(\zeta(-3)\)，并判断所有负奇整数处的符号是否相同。
\end{exercise}
\begin{solution}
代入 \(s=-3\)，得 \(2^{-3}\pi^{-4}\Gamma(4)\zeta(4)=1/120\)。一般地，\(s=1-2m\) 处除正弦外各因子均正，正弦为 \((-1)^m\)，故符号交替。有限的延拓值不是原发散级数的通常和。
\end{solution}
\begin{exercise}\label{b3:ex:26-2}
对整数 \(m\ge1\) 求 \(\zeta'(-2m)\) 与 \(\zeta(2m+1)\) 的关系。
\end{exercise}
\begin{solution}
负偶数处只有正弦因子消失，其导数为 \((\pi/2)(-1)^m\)。因此
\[
 \zeta'(-2m)=\frac{(-1)^m(2m)!}{2(2\pi)^{2m}}\zeta(2m+1).
\]
它非零，另一次核对单零性。
\end{solution}
\begin{exercise}\label{b3:ex:26-3}
证明 \(\xi(1/2+\mathrm{i}t)\) 对实数 \(t\) 为实数且为偶函数。
\end{exercise}
\begin{solution}
共轭将该点送到 \(1/2-\mathrm{i}t\)，函数方程把后者送回原点。因此函数值等于其共轭，同时 \(t\mapsto-t\) 不改变值。
\end{solution}
\begin{exercise}\label{b3:ex:26-4}
计算 \(\eta(1)=\sum_{n\ge1}(-1)^{n-1}/n\)，并用它核验 zeta 在一处的留数。
\end{exercise}
\begin{solution}
在 \(0\le r<1\) 积分几何级数得 \(\sum(-1)^{n-1}r^n/n=\log(1+r)\)，令 \(r\uparrow1\) 并用交错余项控制得 \(\log2\)。又 \(1-2^{1-s}=(s-1)\log2+O((s-1)^2)\)，所以 \(\eta=(1-2^{1-s})\zeta\) 给留数一。
\end{solution}
\begin{exercise}\label{b3:ex:26-5}
证明 \(\int_0^\infty(\vartheta(t)-1)t^{\sigma/2-1}\,\mathrm{d}t\) 对实 \(\sigma\) 收敛当且仅当 \(\sigma>1\)。
\end{exercise}
\begin{solution}
无穷远指数下降。原点由反演有 \(\vartheta(t)-1=t^{-1/2}-1+O(t^{-1/2}e^{-\pi/t})\)，故主项幂次为 \(t^{(\sigma-1)/2-1}\)，可积恰当 \(\sigma>1\)。
\end{solution}
\begin{exercise}\label{b3:ex:26-6}
证明 \(\zeta(\sigma)-1=O(2^{-\sigma})\) 当 \(\sigma\ge2\)。
\end{exercise}
\begin{solution}
提出首项二并积分比较：\(\sum_{n\ge3}n^{-\sigma}\le\int_{5/2}^\infty(t-1/2)^{-\sigma}\,\mathrm{d}t=2^{1-\sigma}/(\sigma-1)\)。所以总和不超过 \(3\cdot2^{-\sigma}\)。更直接地从 \(n=3\) 起用 \(\int_2^\infty t^{-\sigma}\) 也给同一界。
\end{solution}
\begin{optionalexercise}\label{b3:ex:26-7}
证明 \(\sum_\rho|\rho|^{-1-\varepsilon}<\infty\) 对每个 \(\varepsilon>0\) 成立。
\end{optionalexercise}
\begin{solution}
按 \(2^j\le|\rho|<2^{j+1}\) 分层，每层贡献不超过 \(C2^{j+1}(j+1)2^{-j(1+\varepsilon)}\)。该和收敛；有限小半径部分因零点离散且原点非零也有限。
\end{solution}
\begin{optionalexercise}\label{b3:ex:26-8}
说明仅有 \(\xi(s)=\xi(1-s)\) 与实系数为何不足以推出其零点都在临界线，构造一个反例整函数。
\end{optionalexercise}
\begin{solution}
取 \(F(s)=((s-1/2)^2-1)^2+1\)。它为实系数偶中心多项式，具有两种对称性。零点满足 \((s-1/2)^2=1\pm\mathrm{i}\)，其平方根实部非零，因此所有零点都不在 \(\operatorname{Re}s=1/2\)。对称只使零点成组，并不强制它们在对称轴上。
\end{solution}
